\subsection{DECOMPOSABLE ENVELOPE}

The intersection of a family of decomposable Lie subalgebras of \(\mathfrak{gl}(V)\) is clearly decomposable. Consequently, if g is a Lie subalgebra of \(\mathfrak{gl}(V)\) , the set of decomposable Lie subalgebras of \(\mathfrak{gl}(V)\) containing g has a smallest element, called the decomposable envelope of g; in this paragraph, this envelope will be denoted by \(e(\mathfrak{g})\) .

PROPOSITION 4. Let \(\mathfrak{g}\) be a Lie subalgebra of \(\mathfrak{gl}(\mathrm{V})\) and \(\mathfrak{n}\) an ideal of \(\mathfrak{g}\) . Then \(\mathfrak{n}\) and \(e(\mathfrak{n})\) are ideals of \(e(\mathfrak{g})\) , and \([e(\mathfrak{g}), e(\mathfrak{n})] = [\mathfrak{g}, \mathfrak{n}]\) .

Let \(\mathfrak{g}_1\) be the set of \(x\in \mathfrak{gl}(\mathrm{V})\) such that \([x,\mathfrak{n}]\subset [\mathfrak{g},\mathfrak{n}]\) . This is a decomposable Lie subalgebra of \(\mathfrak{gl}(\mathrm{V})\) , containing \(\mathfrak{g}\) and hence \(e(\mathfrak{g})\) , cf.~no. 1, Prop. 3; in other words, \([e(\mathfrak{g}),\mathfrak{n}]\subset [\mathfrak{g},\mathfrak{n}]\) . Let \(\mathfrak{n}_1\) be the set of \(y\in \mathfrak{gl}(\mathrm{V})\) such that

\[
[ e (\mathfrak {g}), y ] \subset [ \mathfrak {g}, \mathfrak {n} ].
\]

This is a decomposable Lie subalgebra of \(\mathfrak{gl}(V)\) containing n by the preceding, and hence containing \(e(\mathfrak{n})\) ; in other words \([e(\mathfrak{g}), e(\mathfrak{n})] \subset [\mathfrak{g}, \mathfrak{n}]\) , so

\[
[ e (\mathfrak {g}), e (\mathfrak {n}) ] = [ \mathfrak {g}, \mathfrak {n} ].
\]

It follows that \([e(\mathfrak{g}),\mathfrak{n}]\subset [e(\mathfrak{g}),e(\mathfrak{n})]\subset \mathfrak{n}\) , so \(\mathfrak{n}\) and \(e(\mathfrak{n})\) are ideals of \(e(\mathfrak{g})\) .

COROLLARY 1. (i) \(\mathcal{D}^i\mathfrak{g} = \mathcal{D}^i e(\mathfrak{g})\) for \(i\geq 1\) , and \(\mathcal{C}^i\mathfrak{g} = \mathcal{C}^i e(\mathfrak{g})\) for \(i\geq 2\) .

\begin{enumerate}
\def\labelenumi{(\roman{enumi})}
\setcounter{enumi}{1}
\tightlist
\item
  If g is commutative (resp. nilpotent, resp. solvable), then \(e(\mathfrak{g})\) is commutative (resp. nilpotent, resp. solvable).
\end{enumerate}

Assertion (i) follows from Prop. 4 by induction on i and (ii) follows from (i).

COROLLARY 2. Let \(\mathfrak{r}\) be the radical of \(\mathfrak{g}\) . If \(\mathfrak{g}\) is decomposable, \(\mathfrak{r}\) is decomposable.

Indeed, \(e(\mathfrak{r})\) is a solvable ideal of \(\mathfrak{g}\) by Prop. 4 and Cor. 1, hence \(e(\mathfrak{r}) = \mathfrak{r}\) .

